\section{Related Work}

\begin{figure*}[t]
    \centering
    \includegraphics[width=0.95\linewidth]{figures/teaser.pdf}
    \caption{We introduce \textbf{OnlineX}, a framework for continuous and progressive 3D scene reconstruction from streaming images. Our core contribution is a active-to-stable state evolution paradigm, which effectively mitigates long-term drift by decoupling the processing of high-fidelity active local details from the maintenance of a stable global structure.}
    \label{fig:teaser}
    \vspace{-12pt}
  \end{figure*}


\paragraph{Generalizable 3D Reconstruction.} 
3D Gaussian Splatting (3DGS)~\cite{kerbl20233d, yu2024mip, ren2024octree, luiten2024dynamic} has recently emerged as a promising alternative for real-time 3D scene reconstruction. While early 3DGS methods require per-scene optimization, generalizable feed-forward models~\cite{chen2024mvsplat, charatan2024pixelsplat, zou2024triplane, zheng2024gps, wewer2024latentsplat, li2024ggrt} were introduced to eliminate this need but remain dependent on pre-computed camera poses from offline SfM tools like COLMAP~\cite{schonberger2016structure}. To overcome this limitation, recent pose-free approaches~\cite{wang2024dust3r, leroy2024grounding, smart2024splatt3r, ye2024no, zhang2025flare} jointly estimate poses and reconstruct scenes. Nonetheless, these methods typically operate on fixed-size sets of images, rather than continuous, online reconstruction. In contrast, our work targets generalizable online 3D reconstruction from streaming RGB input towards real-time applications.

\paragraph{3D Scene Understanding.} 
The integration of 3DGS with vision foundation models like SAM~\cite{kirillov2023segment} and CLIP~\cite{radford2021learning} has recently gained momentum for enabling open-world 3D scene understanding~\cite{qin2024langsplat,ye2024gaussian,xu2024tiger,qu2024goi,ji2025fastlgs,hu2024semantic,qiu2024feature}. Methods like LangSplat~\cite{qin2024langsplat} and Gaussian Grouping~\cite{ye2024gaussian} attach semantic features to Gaussians but are limited by their reliance on per-scene optimization. While some recent feed-forward networks~\cite{li2025scenesplat} offer generalizability, they treat reconstruction and understanding as separate, post-hoc tasks. In contrast, our work introduces a unified, end-to-end framework that learns both visual appearance and language fields concurrently from source images, eliminating the need for per-scene optimization or specialized network modules.

\paragraph{3D Online Paradigm.}
The paradigm of building and interpreting 3D scenes from sequential input is a long-standing goal in computer vision~\cite{engel2014lsd,forster2016svo,teed2021droid,zhu2024nicer,sun2021neuralrecon,zhang2022nerfusion,wang20243d,wang2025continuous,choy20163d,kar2017learning, chen2025long3r}. However, modern learning-based methods often face a fundamental architectural trade-off. Models employing explicit spatial memory like Spann3R~\cite{wang20243d} and LONG3R~\cite{chen2025long3r} incur unsustainable memory overhead, whereas those using a compact, implicit state like CUT3R~\cite{wang2025continuous} are susceptible to geometric drift. Similar trade-offs exist in parallel online perception research~\cite{liu2022ins,huang2021supervoxel,xu2024memory,zhang2020fusion,xia2025scenepainter,wu2023anyview}. Our framework addresses these challenges by introducing a new online paradigm, which leverages the efficiency of implicit representations but resolves their inherent bottleneck by processing active and stable state in two decoupled streams that are then cohesively fused to achieve both high-fidelity detail and long-term consistency.

  % \paragraph{Online 3D Reconstruction and Understanding.}
  % Online reconstruction has been studied in both classical~\cite{engel2014lsd,forster2016svo,teed2021droid,zhu2024nicer} and learning-based methods~\cite{sun2021neuralrecon,zhang2022nerfusion,wang20243d,wang2025continuous,choy20163d,kar2017learning}. SLAM-based pipelines separate tracking and mapping, which is often brittle in real-time. Learning-based models like Spann3R~\cite{wang20243d} and CUT3R~\cite{wang2025continuous} propose spatial memory and state representation but suffer from heavy memory overhead and geometry drift issue when dealing with long-term scene. Concurrently, online scene understanding~\cite{liu2022ins,huang2021supervoxel,xu2024memory,zhang2020fusion} employs memory and feature propagation but often depends on depth inputs and lacks open-vocabulary support. Our approach unifies 3D Gaussian reconstruction and open-vocabulary semantics from RGB streams in a more effective online paradigm.